\section{Considerações Finais}
\label{sec:considerations}

Este estudo combinou análise de regressão e clusterização espacial para investigar como desigualdades socioeconômicas moldam os tempos de acesso a pé em Curitiba. Foi demonstrado que variáveis socioeconômicas como aquelas associadas à escolaridade e empregabilidade explicam variações nos tempos de acesso (Experimento 1), enquanto que a análise de clusterização (Experimento 2) caracterizou grupos intraurbanos heterogêneos levando em consideração aspectos socioeconômicos e tempos médios.

Os resultados contraintuitivos sobre o efeito das variáveis socioeconômicas (g-k) nos tempos de acesso sugerem que a relação entre fatores socioeconômicos e tempo de acesso não é direta. Essa complexidade fica evidente através da análise de clusters, que revela padrões locais que ajudam a explicar essas relações contraditórias. Um exemplo disso são áreas com alta densidade de \ac{pois} que, apesar disso, não apresentam desenvolvimento socioeconômico elevado, indicando que condições socioeconômicas não são os únicos determinantes dos tempos de acesso. No entanto, é importante ressaltar que este estudo não avaliou a qualidade ou a efetividade desses \ac{pois}, o que pode ser determinante para compreender as disparidades observadas.

O estudo não capturou dinâmicas temporais (ex.: variações horárias na oferta de serviços) e focou exclusivamente na acessibilidade temporal a pé, deixando de considerar critérios qualitativos como capacidade de atendimento, condições físicas ou barreiras morfológicas – fatores que influenciam diretamente a efetividade do acesso. Além disso, outros modos de mobilidade, como bicicleta ou transporte coletivo, não foram incluídos na análise, uma limitação a ser explorada em trabalhos futuros.

Parte da redação deste artigo e dos códigos computacionais utilizados foram desenvolvidos com assistência de ferramentas de IA generativa, empregadas para: (i) aprimorar a clareza textual, (ii) sugerir métodos de validação estatística, e (iii) otimizar a visualização de dados. Todo conteúdo gerado foi rigorosamente validado e adaptado, garantindo precisão conceitual e seguindo princípios éticos de pesquisa acadêmica.

\section{Agradecimentos}
\label{sec:acknoledgments}

Ao Ippuc pela disponibilização dos dados de acessos intraurbanos. Ao SocialNet/FAPESP (2023/00148-0) e CNPq (444724/2024-9, 441444/2023-7) pelo apoio. Aos revisores e colegas pelas contribuições.